%%% FIELD proof-point-identification-reductions
Fix (x\in\mathcal X).  We treat the two stated boundary regimes separately.

First suppose that (Gamma_A=1).  By Definition (mathrm{def{:}exposure	ext{-}region}), every ((w_0,w_1)\in D_{A,x}) obeys
\[
w_1(1-w_0)\leq w_0(1-w_1)
\quad\text{and}\quad
w_0(1-w_1)\leq w_1(1-w_0).
\]
The two cross-products are therefore equal.  Expanding and cancelling their common term gives
\[
w_1-w_0w_1=w_0-w_0w_1,
\qquad\text{hence}\qquad
w_1=w_0.
\]
This algebra remains valid at the boundary values (0) and (1).  Write the common value as (w_x).  The prevalence formula in Definition (mathrm{def{:}adjusted	ext{-}arm	ext{-}mean}) then yields
\[
\pi_x=p_xw_x+(1-p_x)w_x=w_x=w_{a,x},
\qquad a\in\{0,1\}.
\]
For every feasible outcome-odds value (	heta>0), Definition (mathrm{def{:}adjusted	ext{-}arm	ext{-}mean}) consequently gives
\[
q_{a,x}(w_x,w_x;\theta)-r_{a,x}
=(\pi_x-w_{a,x})(m_1-m_0)=0.
\]
Thus, by Definition (mathrm{def{:}common	ext{-}u	ext{-}endpoints}),
\[
q^-_{a,x}(w_x,w_x)=q^+_{a,x}(w_x,w_x)=r_{a,x}
\]
and
\[
H^-_x(w_x,w_x)=H^+_x(w_x,w_x)=r_{1,x}-r_{0,x}.
\]
Because every point of (D_{A,x}) is diagonal when (Gamma_A=1), both endpoint objectives are constant on the whole feasible region.  Hence
\[
\ell_x=\upsilon_x=r_{1,x}-r_{0,x}.
\]
Summing over the finite support (mathcal X), as prescribed by Definition (mathrm{def{:}common	ext{-}u	ext{-}endpoints}), proves
\[
L=R=\sum_{x\in\mathcal X}\lambda_x(r_{1,x}-r_{0,x}).
\]
The strictly positive pair (w_{0,x}=w_{1,x}=1/2) belongs to (D_{A,x}).  The exact equality for both the interior and closed maintained classes in Theorem (mathrm{thm{:}sharp	ext{-}common	ext{-}u	ext{-}envelope}) therefore shows that this singleton, rather than merely an outer bound, is the causal ATE identified set and that its endpoint is attained in the interior class.

Now suppose that (Gamma_0=Gamma_1=1).  Then (Gamma_a=1) for each (a\in\{0,1\}).  Take any
\(widetilde P\in\mathcal M^{\circ}(P_{\mathrm{obs}},\Gamma_A,1,1)\), and fix ((a,x)).  Define
\[
c_{u,x}=\widetilde P(U=u,X=x),
\qquad u\in\{0,1\}.
\]
Strict positivity in Definition (mathrm{def{:}positive	ext{-}common	ext{-}u	ext{-}class}) implies
\[
c_{u,x}>0
\quad\text{and}\quad
0<\mu_{a,u,x}<1.
\]
Definitions of (eta) and (mu) give the factorizations
\[
\eta_{a,1,u,x}=c_{u,x}\mu_{a,u,x},
\qquad
\eta_{a,0,u,x}=c_{u,x}(1-\mu_{a,u,x}).
\]
Assumptions (mathrm{ass{:}outcome	ext{-}odds	ext{-}upper}) and (mathrm{ass{:}outcome	ext{-}odds	ext{-}lower}), specialized to (Gamma_a=1), supply the two opposite inequalities and hence the equality
\[
\eta_{a,1,1,x}\eta_{a,0,0,x}
=\eta_{a,1,0,x}\eta_{a,0,1,x}.
\]
Substitution of the preceding factorizations gives
\[
c_{1,x}c_{0,x}\mu_{a,1,x}(1-\mu_{a,0,x})
=c_{0,x}c_{1,x}\mu_{a,0,x}(1-\mu_{a,1,x}).
\]
Division by the strictly positive quantity (c_{0,x}c_{1,x}), followed by expansion and cancellation, yields
\[
\mu_{a,1,x}-\mu_{a,1,x}\mu_{a,0,x}
=\mu_{a,0,x}-\mu_{a,0,x}\mu_{a,1,x},
\qquad\text{so}\qquad
\mu_{a,1,x}=\mu_{a,0,x}.
\]
It remains to identify this common risk.  Assumption (mathrm{ass{:}observed	ext{-}margin}) and the definition of (r_{a,x}) imply
\[
r_{a,x}=\widetilde P(Y=1\mid A=a,X=x).
\]
Assumption (mathrm{ass{:}consistency}) replaces (Y) by (Y(a)) on ({A=a}), while Assumption (mathrm{ass{:}latent	ext{-}ignorability}) removes (A) from the conditional potential-outcome distribution given ((X,U)).  The law of total probability therefore gives
\[
\begin{aligned}
r_{a,x}
&=\widetilde P\{Y(a)=1\mid A=a,X=x\}\\
&=\sum_{u=0}^1
  \widetilde P\{Y(a)=1\mid A=a,U=u,X=x\}
  \widetilde P(U=u\mid A=a,X=x)\\
&=\sum_{u=0}^1
  \mu_{a,u,x}\widetilde P(U=u\mid A=a,X=x)\\
&=\mu_{a,0,x}\sum_{u=0}^1
  \widetilde P(U=u\mid A=a,X=x)
=\mu_{a,0,x}.
\end{aligned}
\]
All displayed conditional probabilities and all divisions used above are well defined because (widetilde P) lies in the strictly positive class; in particular, every ((A=a,U=u,X=x)) cell has positive mass.  We have proved, for every ((a,x)),
\[
\mu_{a,0,x}=\mu_{a,1,x}=r_{a,x}.
\]
Equivalently, (	heta_{a,x}=1) on the strictly positive chart.  This conclusion also follows through the mixture parameterization: Definition (mathrm{def{:}mixture	ext{-}root}) and Lemma (mathrm{lem{:}binary	ext{-}mixture	ext{-}inversion}) give
\[
\kappa(r_{a,x},w_{a,x},1)=r_{a,x},
\qquad m_0=m_1=r_{a,x}.
\]
Since the interval ([Gamma_a^{-1},Gamma_a]) is now the singleton ({1}), Definition (mathrm{def{:}adjusted	ext{-}arm	ext{-}mean}) yields, for every strictly positive feasible posterior pair,
\[
q_{a,x}(w_0,w_1;1)
=r_{a,x}+(\pi_x-w_{a,x})(r_{a,x}-r_{a,x})
=r_{a,x}.
\]
It follows from Definition (mathrm{def{:}common	ext{-}u	ext{-}endpoints}) that
\[
H^-_x(w_0,w_1)=H^+_x(w_0,w_1)=r_{1,x}-r_{0,x}
\]
on the strictly positive part of (D_{A,x}).  Proposition (mathrm{prop{:}closed	ext{-}chart	ext{-}extension}) extends both objectives continuously to all of the compact region (D_{A,x}), including the constant-(U) corners.  Every boundary pair is a limit of strictly positive feasible pairs; hence continuity preserves the displayed constant value on all of (D_{A,x}).  Therefore
\[
\ell_x=\upsilon_x=r_{1,x}-r_{0,x},
\qquad
L=R=\sum_{x\in\mathcal X}\lambda_x(r_{1,x}-r_{0,x}).
\]
Finally, (w_{0,x}=w_{1,x}=1/2) is strictly positive and feasible.  Theorem (mathrm{thm{:}sharp	ext{-}common	ext{-}u	ext{-}envelope}) identifies the images of both (mathcal M^{\circ}(P_{\mathrm{obs}},\Gamma_A,1,1)) and its closed extension with ([L,R]).  Thus the displayed singleton is the exact, interior-attained ATE identified set in the second regime as well.
